\pagenumbering{roman}
\setcounter{page}{1}

\selecthungarian

%----------------------------------------------------------------------------
% Abstract in Hungarian
%----------------------------------------------------------------------------
\chapter*{Kivonat}\addcontentsline{toc}{chapter}{Kivonat}

A szakdolgozat projektem célja egy játék motor készítése, ami a jelenleg is elérhetőkéhez hasonló képességekkel bír.
Mivel ez egy rendkívülien nagy feladat, amit nem lehet a szakdolgozat ideje alatt befejezni, ezért a projekt jó megalapozására fókuszáltam.
Ez azt jelenti, hogy több időt fektettem a megfelelő eszközök kiválasztásába, mint amennyit más esetekben fordítanék erre.
Ilyen volt például a projektben, a kód fordítására használt eszközök megválasztása.
A szakdolgozattal véleményem szerint egy jó alapot tudtam adni a projektemnek, aminek a fejlesztését a jövőben is szeretném folytatni.
Megvalósított részek:

\begin{itemize}
    \item A fordítási környezet megtervezése
    \item Dinamikus, kód betöltés a motor modulokhoz
    \item Kezdeti entitás rendszer
    \item Kód reflexió megvalósítása \Cpp-ban
\end{itemize}

Egy ismerősöm felajánlotta, hogy segít a motor grafikus részének megírásában, azonban a szakdolgozatban csak a saját kódomat fogom bemutatni.
A projekt GitHub-on található és MIT licenszelésű.
A források elérhetőek a következő linken: \url{https://github.com/dpeter99/umbra/tree/thesis} a \verb|thesis| branch-en.


\vfill
\selectenglish


%----------------------------------------------------------------------------
% Abstract in English
%----------------------------------------------------------------------------
\chapter*{Abstract}\addcontentsline{toc}{chapter}{Abstract}

This project's goal is to create a game engine similar in design to currently available big commercial engines.
As this is a large proposal that far exceeds the limits and practicality of a thesis project, I was focusing on giving the project a good foundation.
This means that, I have spent considerable amount of time on investigating solutions to issues other projects would not care about like build systems and project organization.
By the end, I think I have created a solid foundation including the core systems for the engine that I plan to use, and continue working on in the coming years.
Finished parts:

\begin{itemize}
    \item Choosing a build system that is usable with large projects
    \item Dynamic DLL loading of engine modules
    \item Basic Entity system
    \item Reflection system for runtime code metadata
\end{itemize}


A friend of mine has offered to write the renderer for the engine, however, I will be only presenting the parts that are entirely my own.

The project is hosted on GitHub and is MIT licensed.

The source code is available at \url{https://github.com/dpeter99/umbra/tree/thesis} under the \verb|thesis| branch.


\vfill
\selectthesislanguage

\newcounter{romanPage}
\setcounter{romanPage}{\value{page}}
\stepcounter{romanPage}